% !TeX program = lualatex
\documentclass[12pt]{book}

% A4
\usepackage[
  paperheight   = 297mm, 
  paperwidth    = 210mm,
  bindingoffset = 10mm,
  left          = 7.5mm,
  right         = 15mm,
  top           = 20mm,
  bottom        = 22.5mm,
  footskip      = 6mm,
]{geometry}

\usepackage[
  colorlinks,
  allcolors   = black, 
  pdfsubject  = {Go;\ Baduk;\ Weiqi},
  pdftitle    = {Enciclopédia\ de\ Tesujis},
  pdfauthor   = {Philippe\ Fanaro},
  pdfkeywords = {Go;\ Baduk;\ Weiqi;\ Tsumegos},
  pdfproducer = {LaTeX},
  pdfcreator  = {pdflatex;\ bibtex},
  bookmarksnumbered,
]{hyperref}

\def\repoPath{/Users/phili/Code/Go/tsumego_workbooks/src}
\usepackage{\repoPath/config}

\linespread{1.4}

\begin{document}
  \begin{titlepage}
    \setlength\parindent{0ex}
    
    \vspace*{-2.8cm}

    \begin{center}
      \scshape

      \fontsize{38}{38}\selectfont Kifu\\
      \vspace{0.4cm}
      \fontsize{28}{28}\selectfont 棋譜\\
      \vspace{0.4cm}
      \fontsize{23}{23}\selectfont registro de partidas

    \end{center}
    
    \vspace{0.4cm}

    \fontsize{20}{20}\selectfont

    Progredir no Go requer uma reflexão contínua sobre o que fazemos no tabuleiro. No entanto, talvez ainda mais importante do que progredir seja não esquecer.
    
    \begin{figure}[h]
      \centering
      \hbox{
        \hspace{2.2cm}
        \begin{goban}[board dimension       = 12,
                      board size            = 19,
                      scale                 = 1,
                      outline line width    = 0.55mm,
                      horizontal clip start = 1,
                      horizontal clip end   = 19,
                      vertical clip start   = 1,
                      vertical clip end     = 19]
          \parseSgfFile{\repoPath/sgf/kifu/1584_Kashio-Rigen_vs_Honinbo-Sansa.sgf}
        \end{goban}
      }
    \end{figure}

    \vspace{0cm}

    Philippe Fanaro é um jogador amador \emph{dan} brasileiro, formado em Engenharia Elétrica pela Escola Politécnica da USP, que procura trazer mais acesso, em língua portuguesa, ao conteúdo de Go disponível em inglês e línguas asiáticas.
  \end{titlepage}
\end{document}